%! Complex Numbers — Unit 2, Lesson 2.4 — Group Activity (Answer Key)
\documentclass[10pt]{article}
\usepackage{algebra2-article}
\usepackage{algebra2-key}   % pulls in -boxes; do NOT also load -boxes
\newcommand{\TallMath}[1]{$\displaystyle #1\rule[-1.4em]{0pt}{3.2em}$}

\begin{document}

\pageheader{Unit 2, Lesson 2.4}{Group Activity --- Answer Key}

\vspace{0.08in}

\begin{headlinebox}{forestbg}
\small\textbf{The Number $i$.} One new rule, $i^2=-1$, unlocks a whole number system. With your partner,
rewrite negative roots with $i$, put every answer in \textbf{standard form} $a+bi$, and keep the $i^2=-1$
substitution front and center when you multiply.
\end{headlinebox}

\vspace{0.06in}

\begin{tcolorbox}[colback=white, colframe=black!40, title=\textbf{Tier R --- Remediate},
    fonttitle=\bfseries, arc=1.5mm, left=3mm, right=3mm, top=2mm, bottom=2mm, breakable]
\textbf{1. Rewrite with $i$} (simplest radical form).
\[
  \sqrt{-16}=\ans{$4i$} \qquad \sqrt{-49}=\ans{$7i$} \qquad \sqrt{-20}=\ans{$2i\sqrt{5}$}
\]

\textbf{2. Powers of $i$.} Use $i^2=-1$.
\[
  i^2=\ans{$-1$} \qquad i^3=\ans{$-i$} \qquad i^4=\ans{$1$} \qquad i^6=\ans{$-1$}
\]

\textbf{3. Name the parts.} Write $a$ and $b$ for each, and circle \emph{real}, \emph{imaginary}, or \emph{pure imaginary}.
\begin{center}
\renewcommand{\arraystretch}{1.5}
\begin{tabularx}{\linewidth}{@{}l c c X@{}}
\textbf{Number} & \textbf{$a=$} & \textbf{$b=$} & \textbf{Type (circle one)} \\
\hline
$3-5i$ & \ans{$3$} & \ans{$-5$} & real \ / \ \ans{imaginary} \ / \ pure imaginary \\
$-2i$  & \ans{$0$} & \ans{$-2$} & real \ / \ imaginary \ / \ \ans{pure imaginary} \\
$6$    & \ans{$6$} & \ans{$0$} & \ans{real} \ / \ imaginary \ / \ pure imaginary \\
\end{tabularx}
\end{center}

\textbf{4. Add and subtract.} Combine like terms.
\[
  (2+3i)+(4+i)=\ans{$6+4i$} \qquad (7-2i)-(3+5i)=\ans{$4-7i$}
\]
\end{tcolorbox}

\begin{tcolorbox}[colback=white, colframe=black!40, title=\textbf{Tier A --- Approaching Proficiency},
    fonttitle=\bfseries, arc=1.5mm, left=3mm, right=3mm, top=2mm, bottom=2mm, breakable]
\textbf{1. Multiply.} FOIL, then replace $i^2=-1$; give $a+bi$.

\textbf{(a)} $(3+i)(2+i)$
\begin{work}
  (3+i)(2+i) &= 6+3i+2i+i^2 \\
             &= 6+5i+(-1) \\
             &= 5+5i
\end{work}
\textbf{(b)} $(5-2i)(1+4i)$
\begin{work}
  (5-2i)(1+4i) &= 5+20i-2i-8i^2 \\
               &= 5+18i-8(-1) \\
               &= 13+18i
\end{work}
\textbf{(c)} $(4-3i)^2$
\begin{work}
  (4-3i)^2 &= 16-24i+9i^2 \\
           &= 16-24i+9(-1) \\
           &= 7-24i
\end{work}

\textbf{2. Conjugate products.} Multiply each number by its conjugate. (What kind of number do you get?)
\[
  (2+5i)(2-5i)=\ans{$29$} \qquad (1-3i)(1+3i)=\ans{$10$}
\]
\hfill{\small\itshape Always a real number: $(a+bi)(a-bi)=a^2+b^2$.}

\textbf{3. Powers of $i$.} Reduce the exponent mod $4$.
\[
  i^{18}=\ans{$-1$} \qquad i^{25}=\ans{$i$} \qquad i^{44}=\ans{$1$}
\]

\textbf{4. Read the complex plane.} Name each plotted point as $a+bi$.
\begin{center}
\begin{tikzpicture}[scale=0.5]
  \draw[step=1, linegray!40, very thin] (-3,-3) grid (3,4);
  \draw[->, semithick, charcoal] (-3,0)--(3,0) node[below right, font=\scriptsize]{Re};
  \draw[->, semithick, charcoal] (0,-3)--(0,4) node[left, font=\scriptsize]{Im};
  \fill[forest] (2,3) circle (3pt) node[above right, font=\scriptsize]{$P$};
  \fill[forest] (-1,2) circle (3pt) node[above left, font=\scriptsize]{$Q$};
  \fill[forest] (0,-2) circle (3pt) node[right, font=\scriptsize]{$R$};
\end{tikzpicture}
\end{center}
\[ P=\ans{$2+3i$} \qquad Q=\ans{$-1+2i$} \qquad R=\ans{$-2i$} \]
\end{tcolorbox}

\begin{tcolorbox}[colback=white, colframe=black!40, title=\textbf{Tier E --- Extension},
    fonttitle=\bfseries, arc=1.5mm, left=3mm, right=3mm, top=2mm, bottom=2mm, breakable]
\textbf{1. Solve over the complex numbers.} Keep both solutions in $a+bi$ form.
\[
  x^2=-49:\ x=\ans{$\pm7i$} \qquad x^2+9=0:\ x=\ans{$\pm3i$} \qquad (x-2)^2=-16:\ x=\ans{$2\pm4i$}
\]

\textbf{2. Complete the square into complex roots.} Solve $x^2+4x+13=0$.
\begin{work}
  x^2+4x &= -13 \\
  x^2+4x+4 &= -9 \\
  (x+2)^2 &= -9 \\
  x+2 &= \pm3i \\
  x &= -2\pm3i
\end{work}

\textbf{3. Divide (rationalize with the conjugate).} Multiply top and bottom by the denominator's
conjugate, then write in $a+bi$ form.
\begin{work}
  \frac{3+i}{1-2i} &= \frac{(3+i)(1+2i)}{(1-2i)(1+2i)} \\
                   &= \frac{3+6i+i+2i^2}{1+4} \\
                   &= \frac{1+7i}{5} \\
                   &= \frac{1}{5}+\frac{7}{5}i
\end{work}

\vspace{2pt}
\textbf{4. Error analysis.} A student writes $\sqrt{-4}\cdot\sqrt{-9}=\sqrt{36}=6$. Explain the mistake
and give the correct value.
\ansline{The rule $\sqrt{a}\sqrt{b}=\sqrt{ab}$ fails for negatives. Convert to $i$ first: $\sqrt{-4}\cdot\sqrt{-9}=2i\cdot3i=6i^2=-6$.}
\end{tcolorbox}

\end{document}
